\section{Introduction}

Deep Reinforcement Learning (DRL) has significantly transformed the landscape of automated trading systems, offering sophisticated tools for decision-making in complex environments such as the Forex market. This research addresses a critical gap in the application of DRL to Forex trading by proposing a novel integration of the Deep Deterministic Policy Gradient (DDPG) algorithm with Recurrent Neural Networks (RNN), specifically Long-Short Term Memory (LSTM) networks.

\subsection{Motivation}

Automated Forex trading systems traditionally rely on simpler algorithms that often fail to capture the complex and dynamic sequential nature of financial markets. The challenge lies in developing a system that can not only handle the complexities of Forex trading, but also adapt to its dynamic environment in real-time. This research aims to improve automated Forex trading systems by implementing a DDPG algorithm, initially with a Fully Connected Neural Network (FNN) and potentially extending to LSTM networks for better handling of temporal dependencies.

The Forex market, being one of the most liquid and volatile markets, presents unique challenges and opportunities for automated trading systems. A robust DRL-based system that effectively navigates these challenges could significantly improve trading performance and risk management. This research is particularly important as it explores underutilised areas such as the application of DDPG with LSTM networks in Forex trading, and the integration of online learning capabilities through a novel client-server framework with the cTrader \cite{noauthor_forex_nodate} platform. These advances have the potential to set new benchmarks in automated trading efficiency and adaptability.

Developing an effective DRL-based trading system is inherently difficult due to the complex nature of financial markets. Naive approaches often fail due to their inability to handle the high-dimensional and sequential nature of market data, leading to poor generalisation and performance. Additionally, achieving a balance between exploration and exploitation, as well as managing bias-variance trade-offs, are significant challenges that require advanced techniques such as feature engineering, normalisation, regularisation, and adaptive learning rates when optimising the models.

Although there have been significant advances in DRL, the integration of DDPG with LSTM networks specifically for Forex trading remains underexplored. Previous solutions often focus on simpler algorithms or different markets, lacking the sophistication to handle the complexities of the Forex market. Additionally, the implementation of online learning and real-time data processing has been limited by technological constraints and the complexity of developing robust client-server architectures. This research differentiates itself by addressing these gaps, leveraging advanced DRL techniques, and integrating with the cTrader platform to facilitate continuous learning and adaptation.

\subsection{Objectives}

This research proposes a comprehensive trading system comprising several key components:
\begin{itemize}
\item DDPG Algorithm Implementation: The system will initially implement DDPG with an FNN and extend to LSTM networks if feasible to improve the handling of sequential data.
\item Advanced Feature Engineering: To improve model performance, advanced techniques such as feature engineering, normalisation, and technical analysis will be used.
\item Advanced Optimisation Techniques: Techniques like regularisation, e-greedy policies, and Gaussian noise will be used to balance bias/variance and exploration/exploitation problems.
\item cTrader Platform Integration: A novel client-server framework will be developed that will enable the system to process real-time data and execute strategies with online learning capabilities.
\end{itemize}

The proposed timeline from June 2024 to February 2025 ensures a systematic approach to development, testing, and evaluation, with the aim of providing a robust, adaptive, and efficient solution for automated Forex trading. This study is a literature extension of foundational work in DRL, including key contributions such as \cite{mnih_playing_2013}, who demonstrated the potential of deep learning to master Atari games through reinforcement learning with raw pixel input; \cite{mnih_human-level_2015}, who introduced the Deep Q-Network (DQN) algorithm that exceeded human performance in Atari games; \cite{van_hasselt_deep_2016}, who improved DQN by addressing its overestimation bias through the Double DQN (DDQN) algorithm; and \cite{lillicrap_continuous_2015}, who first implemented DDPG, demonstrating its efficacy in solving tasks with continuous action spaces. This research also draws upon the fundamental knowledge presented in key texts such as \cite{goodfellow_deep_2016} on deep learning and \cite{sutton_reinforcement_2018} on reinforcement learning. This research aims to establish a new benchmark in automated Forex trading, advancing the field through innovative applications of DRL technologies and addressing significant existing gaps in the literature.